% Requires: booktabs, tabularx, graphicx.
\chapter{Implementation}
\label{ch:implementation}

This chapter presents the integration of the audio acquisition, display, network and server components of the smart voice note maker. The assembly procedure is based on the firmware-defined connections and the current project wiring documentation.

\section{Step-by-Step Assembly of the Hardware Prototype}
The arrangement uses an STM32 Nucleo-F446RE as the audio and display controller and an ESP32 as the Wi-Fi interface. The INMP441 and OLED are connected to the STM32, while the two controllers exchange commands, audio and status over UART. Each board is intended to receive power through its own USB connection, with a common ground between them. The interconnections are listed in Table~\ref{tab:wiring}.

The hardware integration procedure follows these steps:
\begin{enumerate}
\item Identify the power, ground and signal pins on both development boards and modules.
\item Connect the INMP441 microphone supply, channel selection and I2S signals to the STM32.
\item Connect the SSD1306 supply and I2C signals to the STM32.
\item Cross-connect the STM32 and ESP32 UART transmit and receive signals and establish a common ground.
\item Inspect supply polarity, signal connections and board pin availability.
\item Connect separate USB power to the two boards and program each controller with its respective firmware.
\end{enumerate}

\section{Circuit Integration and Wiring Details}
\subsection{Microphone and Display Connections}
Connect the microphone SCK, WS and SD signals to PB10, PB12 and PC3, respectively. Connect its VDD to 3.3~V, its ground to the common ground and its L/R input to ground for left-channel operation. The project documentation calls for checking the microphone breakout for the recommended 100~k$\Omega$ SD pull-down before adding one. Connect OLED SDA and SCL to PB9 and PB8 and connect its supply and ground after confirming that the module supports 3.3~V. The firmware initializes the display at I2C address \texttt{0x3C}.

\subsection{Controller Interconnection and Power}
Connect STM32 PC6 (transmit) to ESP32 GPIO16 (receive), and ESP32 GPIO17 (transmit) to STM32 PC7 (receive). Connect the grounds together and power the boards independently by USB as specified in the current wiring documentation. Confirm the availability of the ESP32 pins on the actual board and check the wiring and supply polarity before energizing the assembly. The pin assignments follow the current firmware configuration.

\section{Firmware/Code Flashing and Configuration}
\subsection{Build and Configuration}
The STM32 and ESP32 are separate PlatformIO projects in \texttt{firmware/stm32} and \texttt{firmware/esp32}. The former targets \texttt{nucleo\_f446re} and the latter \texttt{esp32dev}, both using the Arduino framework. The STM32 project includes the Adafruit display libraries, and the ESP32 project uses ArduinoJson. The STM32 code additionally configures the I2S and DMA peripherals through the HAL.

For the network client, copy \texttt{firmware/esp32/include/secrets.example.h} to \texttt{secrets.h} in the same directory and configure Wi-Fi, the plain HTTP server origin and the device API key. Configure the server from \texttt{server/.env.example}, including matching device credentials, a session secret, an administrator password and, if AI processing is required, a Gemini API key. Credentials should remain outside the report and version control. The following commands build and upload the two projects independently from the repository root:
\begin{verbatim}
pio run -d firmware/stm32
pio run -d firmware/stm32 -t upload
pio run -d firmware/esp32
pio run -d firmware/esp32 -t upload
\end{verbatim}
The server dependencies are listed in \texttt{server/requirements.txt}; the server is started using \texttt{python run.py} from the \texttt{server} directory after its environment is configured. The default firmware duration is 60~s. If changed, the server's maximum audio duration and command lifetime must be sufficient for the chosen recording and upload.

\subsection{STM32 Audio and Display Firmware}
The STM32 configures SPI2 as an I2S master receiver, using DMA1 stream~3 in circular mode. DMA callbacks copy the selected left-channel data into a 16-entry ring of 128-sample blocks. The foreground code converts each block to a framed 16-bit PCM payload and transmits it over USART6. A full queue or I2S error sets an overflow condition and stops the recording. Display routines initialize the 128~$\times$~64 OLED over I2C and render a title and message. If display initialization fails, messages continue to be sent to the USB serial log; the display is not assumed to be operational.

\subsection{ESP32 Network and Link Firmware}
The ESP32 polls the device-command endpoint, initiates the upload and forwards the server response to the STM32 as a status message. Its UART receive buffer is configured to 32768 bytes. The shared binary protocol uses a two-byte marker, version, type, sequence number, bounded payload length and CRC-16/CCITT-FALSE. The receiver discards malformed frames and the upload path checks audio sequence, size and the final \texttt{END} frame. The ESP32 constructs a 44-byte PCM WAV header and streams the audio in HTTP chunks instead of allocating an entire recording buffer.

\subsection{Server and User Interface}
The FastAPI application stores audio on disk and metadata in a database, SQLite by default. The browser interface allows a logged-in user to request a recording and play saved audio. AI generation is a separate user action; its output and failure state are stored with the note. An optional question-answering endpoint is present on the server, but the split firmware does not call it. Likewise, the older ESP32-only firmware in \texttt{firmware/voice\_note} is a separate implementation and should not be mixed with the split wiring described here.

\section{Challenges and Design Responses}
The following engineering constraints and responses are evidenced by the source and project documentation. They describe design considerations and the corresponding implementation choices.
\begin{enumerate}
\item \textbf{Continuous audio transfer:} A minute of uncompressed PCM occupies approximately 1.92~MB. The design uses DMA blocks, UART frames and chunked HTTP streaming to avoid storing the entire recording on either controller.
\item \textbf{Transfer integrity and buffering:} The finite capture queue can overflow and the serial or network link can lose or delay data. Sequence numbers, CRC checking, timeouts and abort messages allow the firmware to reject incomplete transfers.
\item \textbf{Microphone startup:} The STM32 discards 40 initial audio blocks before signalling readiness, as implemented in the recording routine.
\item \textbf{Pin allocation:} The microphone, display and inter-controller link each require distinct peripheral signals. The design assigns I2S2, I2C and USART6 to separate pins, as listed in Table~\ref{tab:wiring}.
\item \textbf{AI and network dependence:} The STM32 does not run the transcription model. Optional AI generation requires the configured server and external Gemini service; upload also depends on network connectivity. The split firmware does not provide offline recording or automatic retry of a delivered command.
\end{enumerate}

\section{Photographs of the Hardware Model}
The following placeholders should be replaced with actual project photographs after checking that they show the stated configuration.
\begin{figure}[htbp]
\centering
\fbox{\parbox[c][45mm][c]{0.8\textwidth}{\centering Photograph placeholder: overall STM32--ESP32 assembly}}
\caption{Overall assembly of the voice note maker (photograph to be supplied and verified).}
\label{fig:assembly-photo}
\end{figure}
\begin{figure}[htbp]
\centering
\fbox{\parbox[c][45mm][c]{0.8\textwidth}{\centering Photograph placeholder: microphone, OLED and controller wiring}}
\caption{Microphone, display and inter-controller wiring (photograph to be supplied and verified).}
\label{fig:wiring-photo}
\end{figure}
\begin{figure}[htbp]
\centering
\fbox{\parbox[c][45mm][c]{0.8\textwidth}{\centering Screenshot placeholder: recorded note and optional AI output}}
\caption{Web interface showing a recorded note and its AI output (screenshot to be supplied and verified).}
\label{fig:web-photo}
\end{figure}


